\chapter{Clean-Source Assumption \& Fail-Closed Halt Semantics}
\label{ch:legend}

The atlas translates a source whose every symbol it knows. A typo, an idiom and an unknown symbol each lie outside
the charts. This chapter states the invariant that governs them, and \texttt{Lng/any.lean} and
\texttt{Lng/legend.lean} prove each part of it.

\section{The invariant}

\begin{definition}
A chart is fail-closed where it gives no value to a symbol outside it:
\[
  \sigma \notin \Sigma \implies \lBrack\sigma\rBrack(x, y) = \bot \quad \text{for all } x, y \in W_n.
\]
\end{definition}

The structure \texttt{Chart n s} holds this law as its field \texttt{closed}, beside the field \texttt{sound}. Lean
accepts no chart without the proof of both. The charts of \texttt{.g}, C and Python each carry it.

\begin{theorem}[\texttt{halt}, \texttt{no\_lift}]
A symbol outside a chart has no value on any pair of words. It therefore lifts to no symbol that has a value.
\end{theorem}

\begin{definition}
A corpus is a finite list of symbols. Under the clean-source assumption each symbol of a corpus is a symbol of the
chart.
\end{definition}

\begin{theorem}[\texttt{clean\_or\_halt}]
Each symbol of a corpus is a symbol of the chart, or the chart halts on it.
\end{theorem}

The assumption is therefore no hidden premise. Where it fails, the translation stops at the symbol that breaks
it, and it gives no value there.

\begin{theorem}[\texttt{typo}]
For all $x, y \in W_{32}$, the typo \texttt{ad} of add halts in \texttt{.g}, in C and in Python.
\end{theorem}

\section{The canonical chart and the defect legend}

A dirty corpus holds known defects: typos, idioms and specific 1:1 mappings. A defect legend names each of them
and the symbol of the canonical chart that it stands for.

\begin{definition}
Let $C$ be a sound, fail-closed chart. A defect legend $L$ over $C$ is a finite list of entries $(d, t)$, with these
three laws:
\begin{enumerate}
  \item Disjointedness: $d \notin \Sigma_C$ for each entry. Thus $C \cap L = \emptyset$.
  \item Target closure: $t \in \Sigma_C$ for each entry.
  \item One to one: no defect $d$ has two entries.
\end{enumerate}
\end{definition}

The structure \texttt{Legend C} holds the three laws as its fields \texttt{disjoint}, \texttt{target\_closure} and
\texttt{one\_to\_one}. The chart $C + L$ holds the symbols of $C$ and the defects of $L$. It resolves each defect to
the symbol that its entry names, and it gives that symbol's value.

\begin{theorem}[\texttt{extend}]
$C + L$ is a sound, fail-closed chart.
\end{theorem}

\begin{theorem}[\texttt{conservative}, \texttt{faithful}]
On each symbol of $C$, the chart $C + L$ gives what $C$ gives. For each entry $(d, t)$, it gives $d$ the value and
the anchor of $t$.
\end{theorem}

\begin{theorem}[\texttt{defect\_lifts}]
Each defect of $L$ lifts to each symbol of any sound chart with the anchor of its target, at every width and sign.
\end{theorem}

The legend therefore changes nothing on clean source. It adds each known defect with exactly the meaning of the
symbol that it names, and every lift of that symbol holds for the defect too.

\section{The decomposition of a dirty corpus}

Each symbol $\sigma$ of a dirty corpus falls in one of three parts. It is canonical where $\sigma \in \Sigma_C$.
It is a defect where $\sigma$ is a defect of $L$. It is unknown where it is neither.

\begin{theorem}[\texttt{decompose}]
Each symbol is in exactly one of the three parts, and $C + L$ halts on each unknown symbol.
\end{theorem}

\begin{theorem}[\texttt{count}]
The three parts of a corpus together hold as many symbols as the corpus.
\end{theorem}

\begin{definition}
A hole of a corpus is a pair $(p, \sigma)$, where $\sigma$ is the unknown symbol at the place $p$, counted from 0.
\end{definition}

\begin{theorem}[\texttt{holes\_exact}]
$(p, \sigma)$ is a hole that \texttt{holes} reports exactly where the symbol at the place $p$ of the corpus is
$\sigma$, and $\sigma$ is neither a symbol of $C$ nor a defect of $L$.
\end{theorem}

\begin{theorem}[\texttt{translate\_ok}, \texttt{complete}]
\texttt{translate} gives the corpus resolved to symbols of $C$ where the corpus has no hole. Otherwise it gives the
error that lists each hole. Where a corpus has no hole, each of its symbols gives the relation of its anchor.
\end{theorem}

The legend is sound by \texttt{extend}, \texttt{conservative} and \texttt{faithful}. It is complete for a corpus
exactly where \texttt{holes} reports nothing, and \texttt{holes\_exact} makes that test exact in both directions.
An unknown symbol is never translated. It halts, and the error names each hole with its place, for a human to add
to the chart or to the legend.

\begin{proposition}[\texttt{dirty\_holes}, \texttt{dirty\_count}, \texttt{dirty\_translate}, \texttt{idiom}]
Let the legend over C hold the typo \texttt{ad} for add, the idiom \texttt{x = x + y} for add\_assign and
\texttt{xor} for bitxor. The corpus \texttt{add}, \texttt{ad}, \texttt{x = x + y}, \texttt{shl}, \texttt{xor},
\texttt{addd} has two canonical symbols, three defects and one hole, \texttt{addd} at the place 5. Its translation
is the error that names that hole. The idiom gives the value of the Python add\_assign on every pair of words.
\end{proposition}

\section{The C code}

\texttt{lng\_lift} refuses a word that the tables do not hold. Its refusal names the word and the table that lacks
it, as \texttt{the transpiler has no word ad}. The test of \texttt{lng\_check} holds that refusal for the typo
\texttt{ad}. The place of the word in its source is a fact of the caller, and \texttt{lng\_lift} does not report
it. Lean proves \texttt{translate}, and it does not prove the C code.
